\documentclass[12pt, a4paper]{report}
\usepackage{geometry}
\usepackage{xcolor}
\usepackage{sectsty}
\usepackage{fancyhdr}
\usepackage{amsmath}

\geometry{margin=1in}

\chapterfont{\color{red}\huge\bfseries}
\sectionfont{\color{magenta}\large\itshape}
\subsectionfont{\color{orange}}

\pagestyle{fancy}
\fancyhf{}
\lhead{Your Name}
\chead{Introduction to LaTeX}
\rhead{\today}
\cfoot{\thepage}

\title{{\Huge \textbf{Introduction to LaTeX}}}
\author{{\Large Your Name} \\ Your ID}
\date{\today}

\begin{document}

\maketitle
\thispagestyle{empty}

\begin{abstract}
This document is a brief introduction to \LaTeX{}. It covers the report document class,
chapters, custom heading styles, page headers and footers, and other formatting features.
\end{abstract}

\tableofcontents

\chapter{Introduction}
\LaTeX{} is a typesetting system widely used in academia for producing high-quality documents.

\section{Why LaTeX?}
\LaTeX{} is especially useful for documents with complex mathematics, structured formatting,
and consistent styling across large documents.

\subsection{Use Cases}
\begin{itemize}
    \item Research papers
    \item Theses and dissertations
    \item Technical reports
\end{itemize}

\chapter{Mathematics in LaTeX}
\LaTeX{} provides powerful tools for typesetting equations.

\section{Equations}
The quadratic formula is:
\begin{equation}
    x = \frac{-b \pm \sqrt{b^2 - 4ac}}{2a}
\end{equation}

\end{document}
